% ---------------------------------------------------------------------------
% Posterior Distortion from Hypothesis-Dependent Validity Gating in Runtime
% Diagnosis.
%
% This file WAS a build scaffold. It is now the manuscript.
%
% EVERY PROPOSITION IN THIS PAPER IS LABELLED KNOWN. See
% research/math-paper/NOVELTY_AUDIT.md, which is enforced by
% tests/test_math_paper_proofs.py. The contribution claimed is the runtime
% formulation, the static checkability of the condition, and the implemented
% instance -- not the algebra.
%
% Sources:
%   research/math-paper/FORMAL_CORE.md        the proofs
%   research/math-paper/proof_audit.py        symbolic + numeric verification
%   research/math-paper/PRIOR_ART_MATRIX.csv  novelty status of each result
%   research/systems-paper/                   the implemented instance
%
% BIBLIOGRAPHY STATUS: every entry in refs.bib is UNVERIFIED. Egress from the
% authoring environment denies api.crossref.org, doi.org and dblp.org. This is
% a submission blocker and `make bib-guard` reports it. Do not remove a marker
% without checking the publisher record.
% ---------------------------------------------------------------------------
\documentclass[11pt,a4paper]{article}

\usepackage[utf8]{inputenc}
\usepackage[T1]{fontenc}
\usepackage{amsmath,amssymb,amsthm}
\usepackage{booktabs}
\usepackage{graphicx}
\usepackage{longtable}
\usepackage{url}
\usepackage[margin=1in]{geometry}
\usepackage[hidelinks]{hyperref}

\newcommand{\code}[1]{\texttt{\small #1}}
\newtheorem{proposition}{Proposition}
\newtheorem{corollary}[proposition]{Corollary}
\theoremstyle{definition}
\newtheorem{assumption}{Assumption}
\newtheorem{remark}{Remark}

\title{Posterior Distortion from Hypothesis-Dependent Validity Gating\\
in Runtime Diagnosis}

\author{SWORT Aerospace}
\date{}

\begin{document}
\maketitle

\begin{abstract}
A runtime diagnostic layer decides, tick by tick, which of its evidence is
\emph{applicable}. That decision is a coarsening mechanism the system designed
for itself, and whether it may be ignored is a property of the software rather
than of the data. We give the exact algebra. Writing $q$ for the posterior the
system would have obtained from \emph{all} the evidence it could in principle
have scored and $p$ for the posterior it actually produces,
$p_h = q_h w_h / \mathbb{E}_q[w]$, so every pairwise log-posterior ratio
acquires an additive, evidence-independent offset $\log(w_a/w_b)$ that further
data of the same kind cannot argue with. For the additive gated scores runtime
layers compute, $w_h$ is exhibited in closed form and the offset is exactly the
withheld distinguishing contribution. When no active term distinguishes two
hypotheses their ratio is pinned to the prior ratio exactly, for as long as the
gate stays shut, however concentrated the posterior becomes --- a statement
about the produced posterior alone, which does not depend on the choice of
baseline.
We give the condition under which the distortion vanishes, its multivalued
caveat, a sufficient condition for the MAP to be unchanged, and exact
divergence identities:
$\mathrm{KL}(q\Vert p)$ is precisely the Jensen gap of $\log w$ under $q$, so
the distortion is zero exactly when the validity factors do not vary across
hypotheses.

\textbf{We claim no new mathematics.} Each result is an instance of something
known --- self-normalised reweighting, argmax stability under bounded
perturbation, Jensen's inequality --- and the setting is the coarse-data one,
so the paper concedes rather than claims at every step. What we
offer is the runtime-diagnostic formulation, the observation that the
condition is decidable \emph{statically, from a signature table, before the
system runs}, and a deployed layer that fails it: an independently written
checker that imports none of the production analysis agrees with it on
$378{,}000$ of $378{,}000$ tick-by-pair comparisons and finds one hypothesis
pair pinned to its prior for $200$ consecutive ticks.
\end{abstract}

%======================================================================
\section{The problem}

Diagnostic systems that run online do not evaluate every test at every
instant. A test whose preconditions are unmet is \emph{not applicable}, and a
common and natural implementation is to skip it: the corresponding term is
removed from the score. This is not a bug. Scoring a test whose preconditions
are unmet would inject noise, and down-weighting it requires a weight nobody
has measured.

The difficulty is that applicability is often \emph{hypothesis-relative}, and
that the condition producing the fault may be the same condition that removes
the applicability. A receiver that loses its fix stops emitting the observable
that would have named the loss. The diagnostic layer is then reasoning about a
hypothesis using evidence the hypothesis itself destroyed.

Two questions follow, and this paper answers them exactly.

\begin{enumerate}
\item How far can the produced posterior be from the posterior the surviving
      evidence supports, and in what direction?
\item When is the difference nil, and is that condition checkable?
\end{enumerate}

\paragraph{What this paper does not do.} It makes no observability claim: the
physical fault may be perfectly detectable by other means, and in the system we
measure it is. It offers no new inequality, no new estimator and no
calibration. It does not show that any system was unsafe.

%======================================================================
\section{Setting}
\label{sec:setting}

Let $H=\{h_1,\dots,h_n\}$, $n\ge 2$, be a finite set of mutually exclusive
hypotheses, $\pi_h$ the prior, $L_h$ a likelihood under $h$, and $w_h>0$ a
hypothesis-dependent validity factor. Define
\begin{equation}
q_h \;=\; \frac{\pi_h L_h}{\sum_j \pi_j L_j},
\qquad
p_h \;=\; \frac{\pi_h w_h L_h}{\sum_j \pi_j w_j L_j}.
\label{eq:qp}
\end{equation}

\begin{assumption}[strict positivity]\label{a:pos}
$\pi_h, L_h, w_h \in (0,\infty)$ for all $h$.
\end{assumption}

\noindent
\textbf{A closed gate does not violate Assumption~\ref{a:pos}.} The reader has
just met the word \emph{skip}, and it is natural to read a skipped feature as
$w_h=0$. It is not: by \eqref{eq:wgate} a shut gate gives
$w_h=\exp(-\ell_j(h))$, strictly positive and finite for any finite
$\ell_j(h)$. What Assumption~\ref{a:pos} excludes is a zero \emph{likelihood}
--- the \emph{exclusion} of a hypothesis, a different operation with different
semantics, and out of scope here.
\begin{assumption}[finiteness]\label{a:fin}
$n<\infty$.
\end{assumption}
\begin{assumption}[no independence is assumed]\label{a:indep}
Nothing below assumes the gate is independent of the residual, or of the
fault. The gate may be an arbitrary function of the same state the evidence is
a function of.
\end{assumption}
\begin{assumption}[$w$ is exogenous to the algebra]\label{a:w}
$w$ is whatever the implementation makes it. We say what follows \emph{given}
$w$; we do not model how $w$ arises, except in
Section~\ref{sec:ignorable} where a structural condition on it is stated
explicitly.
\end{assumption}

Assumption~\ref{a:indep} is worth dwelling on. A reader familiar with
missing-data arguments will expect an independence or ignorability condition
to be doing the work somewhere. It is not. The identities below are
unconditional; ignorability appears only in Section~\ref{sec:ignorable}, as
the special case in which the distortion vanishes.

\subsection{What $q$ is measured against, and why it matters}
\label{sec:baseline}

Everything below compares two posteriors, so the \emph{baseline} $q$ must be
named. There are two candidates and they are not interchangeable.

\begin{description}
\item[Reading A --- active-evidence baseline.] $L_h$ is the likelihood of the
      evidence that survived the gate. Then the gating is already inside
      $L_h$, there is nothing left for $w$ to carry, $w_h\equiv 1$, and
      $p=q$ identically.
\item[Reading B --- full-evidence baseline.] $L_h$ is the likelihood of
      \emph{all} the evidence the system could in principle have used, and
      $w_h$ carries what the gate withheld. Then $p\ne q$ in general.
\end{description}

\textbf{Reading A makes every result in this paper vacuous, and we say so
before a reader discovers it.} Under Reading A, Propositions
\ref{p:reweight}, \ref{p:logratio}, \ref{p:kl} and \ref{p:bounds} are true and
empty: they describe a reweighting by the constant $1$. An earlier draft of
this paper defined $L_h$ that way and presented Section~\ref{sec:instance} as
an instance of the framework; it was not one.

\textbf{Everything below uses Reading B}, and Proposition~\ref{p:gated} makes
it concrete for the additive scores that runtime diagnostic layers actually
compute --- including the one in Section~\ref{sec:instance}, whose $w$ we
exhibit in closed form.

\begin{remark}[$q$ is counterfactual, and the main claim does not need it]
\label{r:counterfactual}
Under Reading B the baseline $q$ is generally \emph{not computable at runtime},
and in the case Section~\ref{sec:instance} measures it is not computable at
all: the gate closes because the receiver has stopped reporting, so the
withheld $\ell_j(h)$ has no observed value to take. $q$ is a reference point,
not a quantity we claim to evaluate, and $w_h$ is the exact bookkeeping of what
separates the two --- no more than that.

This matters less than it appears, because \textbf{the paper's empirical claim
does not mention $q$}. Proposition~\ref{p:pin} and Corollary~\ref{c:pinned}
relate $p$ to the \emph{prior}, both of which the system computes. A reader who
rejects Reading B entirely still owes an answer to Section~\ref{sec:instance},
where a posterior ratio sat at its prior value for $200$ consecutive ticks.
What such a reader loses is the quantitative apparatus ---
Propositions~\ref{p:kl} and \ref{p:bounds} measure a distance to a baseline,
and a baseline that cannot be evaluated cannot be measured against.
\end{remark}

Write $\mathbb{E}_q[w]=\sum_h q_h w_h$, $w_{\max}=\max_h w_h$,
$w_{\min}=\min_h w_h$, and
\begin{equation}
R \;=\; \log\frac{w_{\max}}{w_{\min}} \;\ge\; 0 ,
\end{equation}
the \emph{log range} of the validity factors.

%======================================================================
\section{The reweighting identity}

\begin{proposition}[reweighting]\label{p:reweight}
Under Assumptions~\ref{a:pos}--\ref{a:fin}, for every $h$,
\begin{equation}
p_h \;=\; q_h\,\frac{w_h}{\mathbb{E}_q[w]} .
\label{eq:p1}
\end{equation}
\end{proposition}

\begin{proof}
Let $Z_q=\sum_j \pi_j L_j$ and $Z_p=\sum_j \pi_j w_j L_j$, both finite and
positive. Then
$q_h w_h/\mathbb{E}_q[w]
= (\pi_h L_h/Z_q)\,w_h \big/ \sum_j (\pi_j L_j / Z_q) w_j
= \pi_h w_h L_h / Z_p = p_h$.
\end{proof}

The normalising constants cancel, and that cancellation is the whole content:
$p$ is the \emph{self-normalised} reweighting of $q$ by $w$.

\begin{proposition}[log-ratio decomposition]\label{p:logratio}
For any $a,b\in H$,
\begin{equation}
\log\frac{p_a}{p_b}
\;=\;
\log\frac{q_a}{q_b}
\;+\;
\log\frac{w_a}{w_b} .
\label{eq:p2}
\end{equation}
\end{proposition}

\begin{proof}
By \eqref{eq:p1}, $p_a/p_b=(q_a/q_b)(w_a/w_b)$; $\mathbb{E}_q[w]$ cancels.
\end{proof}

Equation~\eqref{eq:p2} is the object the rest of the paper is about. The
gating contributes an \textbf{additive, evidence-independent} shift to every
pairwise log-posterior ratio. It does not scale the evidence, does not
interact with it, and does not diminish as evidence accumulates.
$\log(w_a/w_b)$ is a constant offset the data cannot argue with.

%======================================================================
\section{Pinning, and the condition that causes it}
\label{sec:pinning}

\begin{proposition}[pinned to the prior]\label{p:pin}
If $L_a=L_b$ and $w_a=w_b$, then $p_a/p_b=\pi_a/\pi_b$ exactly, for every
value of the observed evidence.
\end{proposition}

\begin{proof}
$q_a/q_b=\pi_a/\pi_b$ since $L_a=L_b$; apply \eqref{eq:p2} with
$\log(w_a/w_b)=0$.
\end{proof}

Proposition~\ref{p:pin} is stronger than "the evidence is weak". The ratio is
not merely slow to move; it \emph{cannot} move. No quantity of further
evidence of this kind changes it.

\paragraph{When does $L_a=L_b$ hold in practice?} For an additive score over
per-feature terms --- the naive-Bayes signature form used by the system in
Section~\ref{sec:instance}, and by many runtime diagnostic layers --- write
$\ell_j(h)$ for feature $j$'s log-likelihood contribution under $h$ and
$g_j\ge 0$ for its gate. Then
\begin{equation}
\log L_a - \log L_b
\;=\;
\sum_{j \in \mathrm{ACTIVE}} \big(\ell_j(a) - \ell_j(b)\big),
\qquad
\mathrm{ACTIVE}=\{j: g_j>0\}.
\end{equation}
Every feature contributing identically to $a$ and $b$ vanishes from this sum
identically, for every value of the observation. Define the
\emph{active distinguishing set}
\begin{equation}
D(a,b) \;=\; \{\, j \in \mathrm{ACTIVE} : \ell_j(a) \neq \ell_j(b) \,\}.
\label{eq:D}
\end{equation}
If $D(a,b)=\emptyset$ then $L_a=L_b$.

\subsection{The gated-score form, which is what runtime layers compute}
\label{sec:gated}

The abstract $w_h$ of Section~\ref{sec:setting} is not something an
implementer recognises. This makes it concrete, and closes the gap
Section~\ref{sec:baseline} opened.

\begin{proposition}[gated score]\label{p:gated}
Let the score be additive over per-feature terms, $\log L_h^{\mathrm{gated}}
= \sum_j g_j\,\ell_j(h)$ with gates $g_j\in[0,1]$, and let $q$ be the
posterior from the ungated score ($g_j\equiv 1$, Reading B). Then $p$ is the
reweighting of $q$ by
\begin{equation}
w_h \;=\; \exp\!\Big(-\sum_j (1-g_j)\,\ell_j(h)\Big),
\label{eq:wgate}
\end{equation}
Proposition~\ref{p:reweight} applies, and for any pair
\begin{equation}
\log\frac{p_a}{p_b} - \log\frac{q_a}{q_b}
\;=\;
-\sum_j (1-g_j)\big(\ell_j(a)-\ell_j(b)\big).
\label{eq:withheld}
\end{equation}
\end{proposition}

\begin{proof}
$\log L_h^{\mathrm{gated}} = \sum_j \ell_j(h) - \sum_j (1-g_j)\ell_j(h)
= \log L_h + \log w_h$ with $w_h$ as in \eqref{eq:wgate}, which is strictly
positive and finite. Apply Proposition~\ref{p:reweight}, then
Proposition~\ref{p:logratio}.
\end{proof}

Equation~\eqref{eq:withheld} says the shift is \textbf{exactly the withheld
distinguishing contribution}: a fully open gate contributes all of its term,
a fully closed one contributes none, and a partial gate contributes in
proportion. And $w_h$ is hypothesis-dependent precisely when some gated
feature contributes differently to two hypotheses --- which is the ordinary
case, not a contrived one.

\begin{corollary}[Proposition~\ref{p:pin}, recovered]\label{c:pinned}
If $g_j=0$ for every $j$ with $\ell_j(a)\ne\ell_j(b)$, then
$\log(p_a/p_b)=\log(\pi_a/\pi_b)$.
\end{corollary}
\begin{proof}
By \eqref{eq:withheld},
$\log(p_a/p_b) = \log(\pi_a/\pi_b) + \sum_j g_j(\ell_j(a)-\ell_j(b))$, and
every surviving summand has $\ell_j(a)=\ell_j(b)$.
\end{proof}

So Proposition~\ref{p:pin} is not a separate result: it is the corner of
Proposition~\ref{p:logratio} in which the gate withholds every distinguishing
term. We state it separately only because it is the case
Section~\ref{sec:instance} measures.

\paragraph{The static-checkability observation.} For a signature model in
which $h$ is characterised by the feature set $\mathrm{LOW}(h)$ it drives low,
$\ell_j(a)\neq\ell_j(b)$ exactly when $j$ lies in the symmetric difference, so
\begin{equation}
D(a,b) \;=\; \big(\mathrm{LOW}(a) \,\triangle\, \mathrm{LOW}(b)\big)
\;\cap\; \mathrm{ACTIVE} .
\label{eq:Dsym}
\end{equation}
The first factor is a property of the \emph{signature table} and can be
computed before the system runs. The second is the gate state. So the question
"can this pair ever be pinned?" is answered by a set operation on a
configuration file, and the question "is it pinned now?" costs one
intersection. This is the observation we regard as the paper's practical
contribution, and Section~\ref{sec:instance} is what it buys.

%======================================================================
\section{Ignorability, and a caveat that is easy to get wrong}
\label{sec:ignorable}

\begin{proposition}[constant validity factor]\label{p:c2}
If $w_h=c$ for all $h$ and some $c>0$, then $p=q$ identically.
\end{proposition}
\begin{proof}
$\mathbb{E}_q[w]=c$, so $p_h=q_h c/c=q_h$ by \eqref{eq:p1}.
\end{proof}

\begin{remark}[what $p$ is, and what ``distortion'' does not
mean]\label{r:marginal}
For binary gates and a score additive over features that the model treats as
conditionally independent given $h$ --- the naive-Bayes form runtime layers
use, and the form assumed by \eqref{eq:wgate} --- the surviving terms multiply
to the marginal likelihood of the surviving features:
\begin{equation}
\prod_{j\in\mathrm{ACTIVE}}\exp\ell_j(h)
\;=\; \mathbb{P}\big(x_{\mathrm{ACTIVE}} \mid h\big),
\qquad\text{so}\qquad
p_h \;\propto\; \pi_h\,\mathbb{P}\big(x_{\mathrm{ACTIVE}} \mid h\big).
\label{eq:marginal}
\end{equation}
\textbf{So $p$ is the posterior given the evidence that survived.} If the gate
were ignorable, $p$ would be the right answer to a smaller question rather
than a wrong answer to the original one, and the gap between $p$ and $q$ would
be information loss rather than error. We state this at the outset because it
is the first objection a missing-data reader will raise and because, under
ignorability, the objection is correct.

It is not the situation of Section~\ref{sec:instance}, for the reason in
Remark~\ref{r:car}. Throughout this paper ``distortion'' names the
multiplicative reweighting of \eqref{eq:p1} and nothing more. The
propositions are identities; whether a particular reweighting is a defect is a
question about the mechanism, and it is answered in the next remark rather
than by the algebra.
\end{remark}

\begin{remark}[the relation to coarsening at random, stated
exactly]\label{r:car}
The gate is a coarsening mechanism, and the framework that owns it is
coarse-data ignorability~\cite{heitjan1991car}; the same setting is examined,
in a stochastic and estimated form, as state-dependent missingness in
latent-state models~\cite{speekenbrink2025statedep}. We claim neither and cite
both.

The relation has to be stated carefully, because it is easy --- and an earlier
draft of this paper did it --- to call Proposition~\ref{p:c2} ``the
ignorability condition''. It is not. Coarsening at random is a condition on
the \emph{mechanism}: the probability of the observed coarsening must not
depend on what was coarsened away. What it licenses is dropping the mechanism
from the likelihood, leaving exactly the $\mathbb{P}(x_{\mathrm{ACTIVE}}\mid
h)$ of \eqref{eq:marginal} --- \emph{not} $p=q$. Proposition~\ref{p:c2} is a
condition on the \emph{withheld evidence}: $w$ constant means nothing
discriminating was withheld, and that does give $p=q$. The second is strictly
stronger, and the two are independent in both directions. A hypothesis-blind
gate that happens to remove a discriminating feature satisfies the first and
not the second; a gate that withholds only non-discriminating features while
being triggered by the fault satisfies the second and not the first.

For the layer of Section~\ref{sec:instance} the mechanism condition fails, and
it fails for the reason the missing-data literature exists to name: the gate
closes \emph{because} the sensor stopped reporting, so the event
$\{\mathrm{ACTIVE}=A\}$ itself carries information about the hypothesis. A
posterior that used that information would carry a factor
$\mathbb{P}(\mathrm{ACTIVE}=A\mid h)$. The implementation carries no such
factor --- a gated feature contributes nothing to the score and the prior is
fixed --- so the produced $p$ is neither $q$ nor the posterior that accounts
for the gate. Supplying the missing factor would require a per-hypothesis
model of gate closure that the system does not have and that we have not
validated; we identify the omission and do not repair it here.

What the paper adds is the reading rather than either condition: because the
coarsening mechanism is a component of the diagnostic software, both questions
above are settled by inspecting a configuration file, before the system runs
and without data.
\end{remark}

\paragraph{The caveat.} The condition is that $w$ be \emph{constant across
hypotheses}, not that the gate be random, and not that the gate be
conditionally independent of the fault. With a binary gate applied uniformly
to all hypotheses, the surviving-evidence factorisation yields a single shared
$w$ and Proposition~\ref{p:c2} applies. With a \emph{multivalued or partial}
gate --- $g_j\in[0,1]$, as real implementations often have --- the per-feature
contribution is $g_j\,\ell_j(h)$, and the hypothesis dependence re-enters
through $\ell_j$ even when $g_j$ itself is hypothesis-blind. Constancy of $w$
then requires the gate-weighted likelihood difference to vanish, which is
Proposition~\ref{p:pin}'s condition and not an independence statement.

This is not hypothetical for the layer in Section~\ref{sec:instance}.
Equation~\eqref{eq:wgate} gives $w_a/w_b=\exp(-S)$ with
$S=\sum_j(\ell_j(a)-\ell_j(b))$ over the withheld features, so $w_a\ne w_b$
exactly when the gate withholds something that separates the pair --- which is
the case we measure. The ignorability condition fails computably, not merely
in principle.

\begin{remark}
\textbf{Conditional independence of the gate from the fault does not by itself
give $p=q$ in the multivalued case.} We state this because the binary
intuition is strong and the multivalued case is the common one.
\end{remark}

%======================================================================
\section{Concentration is not discriminating support}
\label{sec:support}

\begin{corollary}\label{c:support}
If $D(a,b)=\emptyset$ --- equivalently, if the gate withholds every feature
that distinguishes $a$ from $b$ --- then $\log(p_a/p_b)=\log(\pi_a/\pi_b)$
regardless of $\max_h p_h$. In particular $p_a$ may be arbitrarily close to
$1$ --- driven by the prior and by likelihood terms that do not distinguish
$a$ from $b$ --- while the evidence contains nothing separating them.
\end{corollary}
\begin{proof}
Corollary~\ref{c:pinned}; the conclusion involves no property of $\max_h p_h$.
\end{proof}

\begin{remark}[why this is stated through $D$ and not through $w$]
In the abstract setting, pinning is Proposition~\ref{p:pin}: $L_a=L_b$
\emph{and} $w_a=w_b$. Under Reading B that pair of hypotheses is the wrong one
to reach for, because a gate that withholds a distinguishing feature makes
$L_a\ne L_b$ and $w_a\ne w_b$ \emph{both} --- and the ratio is pinned anyway,
because by \eqref{eq:wgate} the two inequalities cancel exactly. The
gate-level statement, Corollary~\ref{c:pinned}, is the one that holds in the
implemented case; Proposition~\ref{p:pin} is its abstract shadow.
\end{remark}

Two different quantities are in play and an implementation that reports only
one of them will mislead its consumer:

\begin{center}
\begin{tabular}{ll}
\toprule
\emph{posterior concentration} & a function of all hypotheses' scores and of
the prior \\
\emph{discriminating support} & a function of which terms are active \\
\bottomrule
\end{tabular}
\end{center}

\begin{remark}[terminology]
We do not call the posterior a \emph{confidence}. A value of $0.9998$ with
$D(a,b)=\emptyset$ carries no information about $a$ versus $b$, and the word
invites a reader to think it does.
\end{remark}

Corollary~\ref{c:support} says what \emph{cannot} be concluded. It does not
say $a$ and $b$ are indistinguishable in principle: another sensor, another
relation, or an ungated version of the same relation may separate them. The
distinction matters practically, because it points at instrumentation rather
than at inference.

%======================================================================
\section{MAP robustness}

\begin{proposition}[sufficient condition]\label{p:map}
Let $\Delta=\log q_{(1)}-\log q_{(2)}$ be the log margin between the top two
hypotheses under $q$. If $\Delta>R$ then
$\arg\max_h p_h=\arg\max_h q_h$.
\end{proposition}

\begin{proof}
By \eqref{eq:p1}, $\log p_h=\log q_h+\log w_h-\log\mathbb{E}_q[w]$, and the
last term is common to all $h$; ranking under $p$ is the ranking of
$\log q_h+\log w_h$. For the $q$-argmax $a$ and any $b\ne a$,
$(\log q_a+\log w_a)-(\log q_b+\log w_b)\ge \Delta-R>0$.
\end{proof}

\begin{remark}[sufficient, not necessary]
Over $20{,}000$ random instances the condition held in $2{,}721$, and in every
one of those the argmax agreed. In $6{,}877$ instances the argmax moved
\emph{while the condition did not hold}, which is permitted. We report the
second number so the proposition is not read as a characterisation.
\end{remark}

Proposition~\ref{p:map} is an instance of the standard argument that an
optimiser is unchanged when the decision margin exceeds the perturbation bound;
we claim it as known.

\begin{remark}[the proposition is silent exactly where the paper is loudest]
\label{r:mapblind}
$R$ is the spread of the withheld log-likelihood across hypotheses, so it is
\emph{largest} in the regime Section~\ref{sec:instance} measures. Sharply: let
$a,b$ be a pinned pair carrying equal prior mass, and suppose they are the top
two under $q$. Writing $S=\sum_j(\ell_j(a)-\ell_j(b))$, \eqref{eq:wgate} gives
$\log w_a-\log w_b=-S$, so $R\ge|S|$, while
$\Delta=|\log(\pi_a/\pi_b)+S|=|S|$. Hence $\Delta\le R$ \emph{always}, with
equality when the pair attains the extremes of $w$: \textbf{Proposition
\ref{p:map} can never certify a pinned pair.} The audit measures this over the
$2{,}282$ sampled instances in which such a pair was the top two and finds a
worst slack of $+8.88\times10^{-16}$, equality to floating-point precision and
no violation.

The equal-prior hypothesis is doing real work and there is an escape, which is
worth naming because it makes the position worse rather than better. Since
$\Delta\le|\log(\pi_a/\pi_b)|+|S|$, the condition $\Delta>R$ forces
$|\log(\pi_a/\pi_b)|>R-|S|\ge 0$: the prior log-odds must carry the entire
excess on their own. So $\Delta>R$ \emph{can} fire on a pinned pair, and when
it does the certificate has been bought with prior mass rather than with
evidence --- a MAP that Corollary~\ref{c:support} has already said carries no
discriminating support.

We state this because it is the honest limit of the result. The MAP condition
is informative when the gate withholds little --- which is when the phenomenon
this paper is about is absent.
\end{remark}

%======================================================================
\section{Exact divergences, and bounds}

\begin{proposition}\label{p:kl}
$\mathrm{KL}(q\Vert p)=\log \mathbb{E}_q[w]-\mathbb{E}_q[\log w]$ and
$\mathrm{KL}(p\Vert q)=\mathbb{E}_p[\log w]-\log \mathbb{E}_q[w]$.
\end{proposition}
\begin{proof}
$\sum_h q_h\log(q_h/p_h)=\sum_h q_h\log(\mathbb{E}_q[w]/w_h)$ by
\eqref{eq:p1}, and symmetrically for the second.
\end{proof}

The first is exactly the Jensen gap of $\log w$ under $q$. It is non-negative,
and zero if and only if $w$ is $q$-almost-surely constant --- which is
Proposition~\ref{p:c2}. So:

\begin{center}
\emph{the gating distorts the posterior precisely to the extent that the
validity factors vary across hypotheses, and by exactly the Jensen gap.}
\end{center}

\begin{proposition}\label{p:bounds}
For every pair $a,b$,
$\big|\log(p_a/p_b)-\log(q_a/q_b)\big|\le R$, and
$\mathrm{TV}(p,q)\le (w_{\max}-w_{\min})/(4\,\mathbb{E}_q[w])$.
\end{proposition}
\begin{proof}
The first is immediate from \eqref{eq:p2}. For the second,
$\mathrm{TV}(p,q)=\tfrac12\sum_h q_h\,|w_h/\mathbb{E}_q[w]-1|
=\mathbb{E}_q|w-\mathbb{E}_q[w]|/(2\mathbb{E}_q[w])$, and the mean absolute
deviation of a variable supported on $[w_{\min},w_{\max}]$ is at most
$(w_{\max}-w_{\min})/2$.
\end{proof}

The first bound is attained; the audit's worst slack over $20{,}000$ instances
is $+1.78\times10^{-15}$, equality to floating-point precision.

\textbf{The second is attained too, and we can exhibit the witness.} Take
$q=(\tfrac12,\tfrac12)$ and a two-point $w=(a,b)$. Then
$\mathbb{E}_q[w]=(a+b)/2$, $p=(a,b)/(a+b)$, and
\begin{equation}
\mathrm{TV}(p,q)=\frac{|a-b|}{2(a+b)}
=\frac{w_{\max}-w_{\min}}{4\,\mathbb{E}_q[w]},
\label{eq:tvattained}
\end{equation}
an identity rather than an inequality, on both ordering branches. So the bound
\textbf{cannot be improved as a function of $(w_{\max},w_{\min},
\mathbb{E}_q[w]$)}. An earlier draft declined to claim this; the construction
is elementary and we should have looked for it.

We still claim no more than that. Attainment on this witness says the constant
$4$ and the functional form are best possible \emph{given those three
arguments}; it says nothing about bounds using more information about $w$, and
it is not a claim of priority --- bounded-likelihood-ratio divergence bounds
are an established subject~\cite{rukhin1997bounded}.

%======================================================================
\section{Verification}

Every proposition above is checked two ways, and the artifact is committed.
Symbolically with a computer-algebra system at $n=4$
(\code{simplify(lhs$-$rhs)$\,=\,$0}, exact), and numerically over $20{,}000$
random instances. Worst absolute deviations:

\begin{table}[h]
\centering
\small
\begin{tabular}{llr}
\toprule
Result & Symbolic ($n=4$) & Numeric worst deviation \\
\midrule
Prop.~\ref{p:reweight} reweighting          & exact & $4.44\times10^{-16}$ \\
Prop.~\ref{p:logratio} log-ratio            & exact & $1.78\times10^{-15}$ \\
Prop.~\ref{p:pin} pinned to prior           & exact & $8.88\times10^{-16}$ \\
Prop.~\ref{p:gated} gated score is a reweighting & --- & $5.55\times10^{-16}$ \\
Prop.~\ref{p:gated} shift $=$ withheld contribution & --- & $3.55\times10^{-15}$ \\
Cor.~\ref{c:pinned} pinned when all distinguishing gated & --- & $1.78\times10^{-15}$ \\
Prop.~\ref{p:c2} constant $w$               & exact & --- \\
Prop.~\ref{p:kl} $\mathrm{KL}(q\Vert p)$    & exact & $1.33\times10^{-15}$ \\
Prop.~\ref{p:kl} $\mathrm{KL}(p\Vert q)$    & exact & $1.53\times10^{-15}$ \\
Prop.~\ref{p:bounds} $R$ bound              & ---   & 0 violations, slack $+1.78\times10^{-15}$ \\
Prop.~\ref{p:bounds} TV bound               & ---   & 0 violations, slack $-8.38\times10^{-5}$ \\
Prop.~\ref{p:map} MAP condition             & ---   & held in $2{,}721/2{,}721$ \\
Rem.~\ref{r:mapblind} pinned pair never certifiable & --- & 0 violations, slack $+8.88\times10^{-16}$ \\
\bottomrule
\end{tabular}
\caption{Verification. The symbolic pass is at $n=4$; the hand proofs are for
general $n$ and the machine check is not.}
\label{tab:verify}
\end{table}

%======================================================================
\section{An implemented instance}
\label{sec:instance}

The condition in \eqref{eq:Dsym} is not hypothetical. A deployed cross-sensor
diagnostic layer for small uncrewed aircraft maintains a naive-Bayes diagnosis
over eight mutually exclusive causes, with a per-feature gate. For the pair
$(\texttt{nominal},\texttt{gps\_fault})$ the symmetric difference is a single
feature: the GNSS internal-physics consistency check. Its gate closes when the
receiver reports no valid fix. So a fault that removes the fix removes the only
term that could name it, and \eqref{eq:Dsym} is empty for the duration.

To test that reading without circularity, a checker was written that evaluates
\eqref{eq:Dsym} directly from the committed signature table and a per-tick
record of which features were active. It imports none of the production
analysis --- a test runs it in a fresh interpreter and asserts that no module
of the diagnostic layer, and no part of its support analysis, appears in
\code{sys.modules} afterwards. Its input carries no leader, no posterior and no
support label. Against the production computation of the same quantity, over
nine scenarios, $1500$ ticks each and all $28$ unordered hypothesis pairs:

\paragraph{What the agreement does and does not establish.} It rules out that
the collapse is an artefact of the production analysis pipeline: two
implementations sharing no code reach the same verdict every time. It does
\textbf{not} validate the signature table. Both evaluate
$\mathrm{LOW}(a)\triangle\mathrm{LOW}(b)$ from the same committed
configuration, so a table that misdescribes what the features measure would be
misdescribed identically $378{,}000$ times. Whether $\mathrm{LOW}(h)$ is
faithful to the sensors is a separate question that no quantity of agreement
can answer, and the structural argument above is downstream of it.

\begin{table}[h]
\centering
\small
\begin{tabular}{lr}
\toprule
Tick-by-pair comparisons & $378{,}000$ \\
Agreements & $378{,}000$ \\
Disagreements & $0$ \\
Cases in the one known disagreement class & $0$ \\
\bottomrule
\end{tabular}
\caption{Independent checker against production. The two definitions differ in
exactly one degenerate case, when a feature's consistency is exactly $0.5$ and
a structurally distinguishing term contributes zero numerically. That class was
predicted before the comparison, counted separately, and did not occur.}
\label{tab:checker}
\end{table}

What it found:

\begin{itemize}
\item $(\texttt{nominal},\texttt{gps\_fault})$ is the \emph{only} pair among
      the $28$ that ever collapses.
\item It collapses for $76$ ticks at the start of \emph{every} scenario,
      including fault-free flight, because the gate is deliberately closed
      during a smoother warm-up. A structural property, not a fault.
\item Under injected fix loss it collapses for a further $200$ consecutive
      ticks --- exactly the injected window --- during which
      Corollary~\ref{c:pinned} says the weight ratio equals the prior
      ratio. (Corollary~\ref{c:pinned} rather than Proposition~\ref{p:pin}:
      the layer's gate is a skip, so it withholds a term that contributes
      differently to the two hypotheses, and the pinning arises from the
      exact cancellation in \eqref{eq:withheld} rather than from $L_a=L_b$.)
\end{itemize}

\paragraph{A negative control on a current stack.} The same predicate was run
over $15{,}595$ ticks from $20$ trials on a current-generation autopilot in
software-in-the-loop, in a campaign whose validity criterion was that the fix
never dropped. Under our account the predicate must then report the pair
distinguishable at every tick, and a single collapse would falsify it. It
reports $0$ collapses in $15{,}595$ ticks. This is a control, not a
replication of the fix-loss behaviour, and we do not present it as one: no
fix-loss arm exists on that platform in our data.

\paragraph{The implemented consequence.} The layer now reports a support state
separately from the normalized hypothesis weight it emits. When
$D(a,b)=\emptyset$ against a required
competitor, the tick is emitted as \emph{support lost for required
information} rather than as eligibility --- \textbf{at any weight, because no
weight threshold enters the rule at all}. That is deliberate: conditioning
identifiability on the weight would make it a function of the very score whose
support is in question. The rule is likewise independent of
freshness, so removing the freshness rule would not rescue it.

The absence of a threshold is visible in the data rather than only in the
code. Across the ticks where the condition fires, the leader's weight ranges
from $0.9691$ to $0.9998$; a rule gated at, say, $0.999$ would have
passed most of them.

%======================================================================
\section{What is new, stated plainly}

Nothing in Sections 3--8 is new mathematics, and we do not want a reader to
have to work that out. Table~\ref{tab:novelty} is the audit.

\begin{table}[h]
\centering
\small
\begin{tabular}{llp{0.42\textwidth}}
\toprule
Result & Status & Prior art \\
\midrule
Prop.~\ref{p:reweight}, \ref{p:logratio} & KNOWN & self-normalised importance
weighting \\
Prop.~\ref{p:pin}, Cor.~\ref{c:support} & KNOWN as algebra & --- \\
Prop.~\ref{p:gated}, Cor.~\ref{c:pinned} & KNOWN & reparameterisation of a
log-linear score; one line of algebra \\
Prop.~\ref{p:c2} & KNOWN as algebra & the zero-loss corner of
Prop.~\ref{p:reweight}. It is \emph{not} the coarsening-at-random
condition~\cite{heitjan1991car}; the two are compared in
Remark~\ref{r:car} \\
Prop.~\ref{p:map} & KNOWN & argmax stability under bounded perturbation \\
Prop.~\ref{p:kl}, \ref{p:bounds} & KNOWN & Jensen; bounded-support mean
absolute deviation \\
\bottomrule
\end{tabular}
\caption{Novelty audit. Every proposition is labelled KNOWN.}
\label{tab:novelty}
\end{table}

What we claim is: the exact runtime-diagnostic formulation; the connection
between hypothesis-dependent validity and posterior distortion stated as
\eqref{eq:p2}; the structural-support reading and its
\emph{static checkability} from a signature table; and the implemented
counterexample of Section~\ref{sec:instance}.

We claim no new information-theoretic result, no new Bayesian identity, and no
sharper bound than what is already available.

%======================================================================
\section{Related work}

The mechanism by which a fault removes its own observability is studied in
discrete-event systems under state-dependent observation
maps~\cite{dong2024statedependent} and, for observation loss imposed from
outside, under intermittent and permanent
loss~\cite{carvalho2012intermittent,carvalho2013permanent} and general sensor
failure~\cite{takai2021diagnosability}. We assert no priority over the general
mechanism. Our setting differs in being a continuous-signal Bayesian layer in
which applicability is hypothesis-relative and decided by the diagnostic system
itself, so the condition is a property of its configuration rather than of an
automaton.

The ignorability question is the coarsening-at-random
condition~\cite{heitjan1991car}, and the same question for a mechanism that
depends on the latent state --- ours does --- is treated in latent-state
missingness models, where the mechanism is a stochastic submodel estimated
from data~\cite{speekenbrink2025statedep}. The difference we can state is that
ours is deterministic and readable from the diagnostic system's own
configuration, so it is decided rather than estimated; we assert no more than
that, and in particular no claim that an estimated mechanism could not
represent it. Which faults a given instrumentation can
separate is the subject of structural isolability and sensor
placement~\cite{krysander2008sensorplacement}; that literature answers the
question our Section~\ref{sec:support} raises but does not answer, namely what
to instrument.

%======================================================================
\section{Limitations}

\begin{enumerate}
\item \textbf{The baseline $q$ is counterfactual and, in the measured case, not
      evaluable} (Remark~\ref{r:counterfactual}). Propositions~\ref{p:kl} and
      \ref{p:bounds} therefore quantify a distance to a reference point rather
      than to anything the system computes. The pinning result does not depend
      on this and is stated in terms of $p$ and the prior alone.
\item \textbf{Proposition~\ref{p:map} is vacuous in the regime we measure}
      (Remark~\ref{r:mapblind}): a pinned pair with equal prior mass satisfies
      $\Delta\le R$ always, so the sufficient condition cannot fire. We have
      no MAP-stability result for the case the paper is about.
\item The symbolic verification is at $n=4$. The hand proofs are general; the
      machine check is not. Propositions~\ref{p:gated} and \ref{p:map} and
      Remark~\ref{r:mapblind} are checked numerically only.
\item The implemented instance is $n=1$: one deployed layer, whose skip-style
      gate is a natural implementation choice but whose prevalence we do not
      establish. A structural-analysis FDI design would mark the fault set
      unisolable instead, and would not exhibit this.
\item The quantities the layer reports are model states, not calibrated
      probabilities, and we do not calibrate them. The results are invariant to
      the layer's tempering constant, which scales $\log q$ for every
      hypothesis alike.
\item No claim of operational harm is made, and none is supported by our data.
\item \textbf{Every bibliographic entry in this manuscript is unverified.} The
      authoring environment could not reach a publisher record. This is a
      submission blocker and is reported as one.
\end{enumerate}

%======================================================================
\section*{A separate body of work, not in this paper}

This project also holds results on \emph{anytime-valid} evidence --- e-value
calibration, blockwise e-processes and mission-level error control --- with
their own register at \code{docs/math/PROOF\_STATUS.md}. That register contains
six refutations, several of them counterexamples to constructions this project
itself proposed, and four open items. It is a different body of work from the
one above and is deliberately \textbf{not} part of this manuscript. We name it
here so that a reader who encounters the register does not mistake its open
items for gaps in the results proved above.

\bibliographystyle{plain}
\bibliography{references}

\end{document}
